\section{Related work}
MMLU supplies diverse subject-level assessments \citep{hendrycks2021mmlu}; HELM situates accuracy within a broader collection of scenarios and metrics \citep{liang2023helm}. BBH illustrates that a benchmark family can itself aggregate heterogeneous tasks \citep{suzgun2022bbh}. Our equal-subject MMLU design and six-column Open LLM design analyze particular task units, rather than reproducing every published leaderboard weighting or evaluating all dimensions of model quality. The Benchmark Lottery shows how benchmark choices affect method comparisons \citep{dehghani2021lottery}, motivating our composition checks.

\citet{lanctot2023socialchoice} already treat tasks as voters and agents as alternatives and study voting methods for evaluation, including Borda. Our comparison uses one ordinal rule alongside four cardinal summaries; it does not survey social-choice methods or test their recommendations for maximal lotteries. Constructive voting bribery asks whether constrained edits can make a preferred candidate win \citep{faliszewski2016bribery}. Target-only upward rank changes offer a useful analogy, but our restricted score-ceiling menu has its own optimization structure; no general bribery-complexity claim is imported.

Contamination studies demonstrate that benchmark-related training information can inflate scores and that cross-task effects matter \citep{zhou2023cheater,yang2023rephrased}. They motivate selective improvement without calibrating our gain menu to a training dose. Selective disclosure and uneven data access in Arena are distinct strategic mechanisms \citep{singh2025illusion}. A designer--developer game with endogenous strategic post-training \citep{chen2026incentives} studies interactions among developers and the benchmark designer. We instead hold rival score rows fixed and ask an exact reachability question for one target. We claim a reproducible conditional comparison of intervention breadth, not the introduction of voting-based evaluation, a prediction of strategic equilibria, evidence that listed models are contaminated, or priority for the general research question.
